\subsection{Определение параметров выборки}

Рассмотрим выборку \( \vec{X}=(X_1,\ldots,X_n)^T\sim \textrm{B}_\theta \), где \( n = 50 \):

\begin{table}[h!]
\centering
\begin{tabular}{|c|c||c|c||c|c||c|c||c|c|}
\hline
№ & Значение & № & Значение & № & Значение & № & Значение & № & Значение \\ \hline \hline
1  & 0 & 11 & 1 & 21 & 0 & 31 & 0 & 41 & 1 \\ \hline
2  & 0 & 12 & 1 & 22 & 1 & 32 & 1 & 42 & 0 \\ \hline
3  & 0 & 13 & 0 & 23 & 1 & 33 & 0 & 43 & 1 \\ \hline
4  & 1 & 14 & 1 & 24 & 1 & 34 & 0 & 44 & 1 \\ \hline
5  & 1 & 15 & 1 & 25 & 1 & 35 & 0 & 45 & 1 \\ \hline
6  & 0 & 16 & 0 & 26 & 1 & 36 & 0 & 46 & 0 \\ \hline
7  & 0 & 17 & 0 & 27 & 1 & 37 & 0 & 47 & 0 \\ \hline
8  & 1 & 18 & 1 & 28 & 0 & 38 & 1 & 48 & 1 \\ \hline
9  & 1 & 19 & 0 & 29 & 0 & 39 & 1 & 49 & 0 \\ \hline
10 & 1 & 20 & 1 & 30 & 1 & 40 & 1 & 50 & 1 \\ \hline
\end{tabular}
\caption{Выборка из распределения Бернулли}
\end{table}

Можно заметить, что закон распределения зависит от некоторой величины~\( \theta \). На самом деле мы указали класс распределений, который целиком определяется значением скаляра \( \theta \). Это и есть \textbf{параметр} \( \theta \), который может принимать значения из множества \( \Theta \):
\[
    \Theta=\{\theta\mid \theta\in(0,1)\}\text{ --- для }\textrm{B}_{\theta}
\]

Мы не можем найти истинное значение параметра \( \theta \), но по данной выборке возможно предоставить оценку \( \hat{\theta} \) при помощи \textit{метода максимального правдоподобия}.